% Shared basic results; load once from the master.
% Dependency order: certificates -> converse -> regularity -> convergence.
\subsection{Positive variational equations and converse contraction weights}
\label{s:PositiveWeights}
Let $\varphi_t$ be the flow of $\dot y=g(y)$ on a compact forward-invariant
set $K$. For a continuous matrix field $A$, write $\Psi(t,y)$ for the
fundamental matrix of
\[
 \dot\xi=A(\varphi_t(y))\xi,\qquad \Psi(0,y)=I.
\]
For a $C^1$ field $H$, put $D_hH=(\partial_yH)g$, componentwise.
The matrix $A$ must be specified: estimates for the full, transverse,
and exterior-power variational equations have different consequences.

\beP[Common and transported certificates for linear equations along a flow]
[][VarCert]
Suppose $g$ is $C^1$ near $K$ and $A:K\to\mathbb R^{d\times d}$ is
continuous. Each of the following conditions implies uniform exponential decay.
\begin{enumerate}
\item A symmetric $C^1$ field $P$ satisfies, for $0<m\le M$ and $c>0$,
\begin{equation}\label{HG:transportLMI}
 mI\preceq P\preceq MI,\qquad
 D_hP+A^\top P+PA\preceq-2cP\qquad(y\in K).
\end{equation}
Then $\|\Psi(t,y)\|_2\le\sqrt{M/m}\,e^{-ct}$.
\item A fixed diagonal signature matrix $\Sigma$, with entries $\pm1$,
makes $C=\Sigma A\Sigma^{-1}$ Metzler on $K$, and a $C^1$ field $v$ satisfies
\begin{equation}\label{HG:transportLP}
 m\mathbf1\le v\le M\mathbf1,\qquad
 D_hv+C^\top v\le-cv\qquad(y\in K).
\end{equation}
Then $\|\Psi(t,y)\|_1\le(M/m)e^{-ct}$.
This assertion also holds for a continuous flow and continuous $v$
with a continuous forward orbital derivative, used in place of $D_hv$.
\end{enumerate}
The norms $(\xi^\top P(y)\xi)^{1/2}$ and $v(y)^\top|\Sigma\xi|$
respectively decay at rate $c$ along trajectories.

For constant certificates the derivative terms vanish. If
$A(K)\subseteq\operatorname{conv}\{A_1,\ldots,A_q\}$, it suffices to check
one of the following systems with the same certificate at every vertex:
\begin{equation}\label{HG:LMI}
 P\succ0,\qquad A_j^\top P+PA_j\preceq-2cP\quad(1\le j\le q),
\end{equation}
\begin{equation}\label{HG:LP}
 v\gg0,\qquad C_j^\top v\le-cv\quad(1\le j\le q),
 \qquad C_j=\Sigma A_j\Sigma^{-1}\text{ Metzler}.
\end{equation}
For fixed $c$, these are LMIs in $P$ and linear inequalities in $v$,
respectively. They express logarithmic-norm bounds in the constant norms
with transformations $\Theta=P^{1/2}$ and
$\Theta=\operatorname{Diag}(v)\Sigma$.
\eeP
\begin{proof}
Along a solution,
\[
 \frac{d}{dt}(\xi^\top P\xi)
 =\xi^\top(D_hP+A^\top P+PA)\xi
 \le-2c\,\xi^\top P\xi,
\]
which proves the first estimate. For the second, positivity of the
Metzler equation for $\Sigma\xi$ gives, almost everywhere,
\[
 \frac{d}{dt}\bigl(v^\top|\Sigma\xi|\bigr)
 \le(D_hv+C^\top v)^\top|\Sigma\xi|
 \le-cv^\top|\Sigma\xi|.
\]
Integration and the bounds on $v$ prove the claim, including its
orbital-derivative version. Convex combinations preserve the constant
inequalities and the Metzler property. The logarithmic-norm statements
follow from the symmetric part in the Euclidean case and the weighted
column sums in the $\ell^1$ case.
\end{proof}

\beD[Lur'e systems are absolutely contracting when one norm contracts them for a whole class of nonlinearities]
[absolute Lur'e contractivity][LurCon]
A \emph{Lur'e system} is $\dot x=-Vx+B\Phi(Wx,t)$, with $V\in\mathbb R^{n_S\times n_S}$,
$B\in\mathbb R_+^{n_S\times m}$, $W\in\mathbb R_+^{m\times n_S}$. Let $P\succ0$, $\|x\|_P^2=x^\top Px$ and
$\eta>0$. The right-hand side $f$ is \emph{strongly infinitesimally contracting} with respect to
$\|\cdot\|_P$ with rate $\eta$ if
\[
 (f(x_1,t)-f(x_2,t))^\top P(x_1-x_2)\le-\eta\|x_1-x_2\|_P^2
 \qquad\text{for all }x_1,x_2,t.
\]
The system is \emph{absolutely contracting} over a class of nonlinearities if this holds with
the same $P$ and $\eta$ for every $\Phi$ in the class. The class of
\cite{davydov2024exponential} consists of the \emph{cocoercive} maps,
\[
 (\Phi(y_1,t)-\Phi(y_2,t))^\top(y_1-y_2)\ge\rho\|\Phi(y_1,t)-\Phi(y_2,t)\|_2^2,
 \qquad\rho>0.
\]
\eeD

\beR[Absolute Lur'e contractivity and the test $\bar R_0<1/L$]
By \cite[Theorem~1]{davydov2024exponential}, absolute contractivity over the cocoercive class
holds if and only if some $\lambda\ge0$ satisfies
\[
 \begin{pmatrix}
 -V^\top P-PV+2\eta P & PB+\lambda W^\top\\
 B^\top P+\lambda W & -2\lambda\rho I
 \end{pmatrix}\preceq0.
\]
A coordinatewise response with $0\le \Phi'\le L$ is cocoercive with $\rho=1/L$, so
$\dot x=-Vx+\Phi(\beta Wx)$ is the case $B=I$, with $\beta W$ in place of $W$. Put
$\bar R_0=\beta\rho(V^{-1}W)$. For $V=I$, $\beta>0$ and symmetric $W\ge0$ with $r=\rho(W)>0$, the
matrix inequality with $P=I$, $\lambda=1/(\beta r)$ reduces to $\beta Lr\le1-\eta$, so it agrees
with the test $\bar R_0<1/L$.

In general the cocoercive class is strictly larger and the two conditions differ, as the
following illustration, which is not taken from \cite{davydov2024exponential}, shows. Take $V=I$,
$W=\begin{pmatrix}0&a\\0&0\end{pmatrix}$ with $\beta aL>2$. Then $\bar R_0=0$, and every coordinatewise
response gives a triangular system, contracting in the norm $\max\{|z_1|/p_1,|z_2|\}$ for any
$p_1>\beta aL$. The linear map
$\Phi(y)=\frac L2\begin{pmatrix}1&-1\\1&1\end{pmatrix}y$ is cocoercive with $\rho=1/L$, but
$-I+\beta\frac L2\begin{pmatrix}1&-1\\1&1\end{pmatrix}W$ has the eigenvalue $-1+\beta aL/2>0$, so no
$P$ gives absolute contractivity.

With $V$ a nonsingular $M$-matrix and $W\ge0$, absolute contractivity over the cocoercive class
implies $\bar R_0<1/L$: the admissible map $\Phi(y)=Ly$ requires $-V+\beta LW$ to be Hurwitz. The
example above shows that the converse fails. Thus the distinction is not merely the choice of
norm: the cocoercive class also allows coupling between the coordinates of the response.
For firing-rate networks with monotone nonexpansive coordinatewise activations,
\cite{Gokhale2026} give a sufficient condition with a positive diagonal multiplier. An elementary
deduction from it recovers the spectral bound exactly for $V=I$, including nonsymmetric and
reducible $W\ge0$. Put $M=\beta LW$ and normalize the response to have slopes in $[0,1]$. The
criterion is
\[
 \begin{pmatrix}
 -2(1-\eta)P & P+M^\top Q\\
 P+QK & -2Q
 \end{pmatrix}\preceq0,
 \qquad P\succ0,\quad Q\succ0\ \text{diagonal},\quad\eta>0.
\]
Testing the identity activation shows that its feasibility implies $\rho(M)<1$. Conversely, if
$\rho(M)<1$, the nonnegative matrix $G=(I+M)/2$ is Schur stable and admits a positive diagonal
matrix $P$ with $G^\top PG\prec P$ \cite{rantzer2015scalable}. Taking $Q=P$, the Schur complement is
\[
 2\bigl(G^\top PG-(1-\eta)P\bigr),
\]
which is negative definite for sufficiently small $\eta>0$. The triangular example above satisfies
this criterion: it separates the coupled cocoercive class from the coordinatewise class, not
the test $\bar R_0<1/L$ from the coordinatewise criterion.
These LMIs use a quadratic function of the difference of two states. The classical
Lur'e--Postnikov quadratic-plus-integral construction, associated with Popov criteria
\cite{Khalil2002}, instead works about an equilibrium. For $\dot x=-Vx+\Phi(Wx+h)$ with $V$ positive
diagonal, slopes in $[0,1]$ and an equilibrium $x^*$, put $e=x-x^*$ and
$\psi_i(s)=\Phi_i((Wx^*+h)_i+s)-\Phi_i((Wx^*+h)_i)$, so that $\dot e=-Ve+\psi(We)$. With $X\succ0$
and $N=\operatorname{Diag}(\nu_i)\succeq0$, $T\succeq0$, both diagonal, the function
\[
 L_{\mathrm{LP}}(e)=e^\top Xe+2\sum_i\nu_i\int_0^{(We)_i}\psi_i(s)\,ds
\]
is bounded above and below by quadratics. Writing $z=\psi(We)$, differentiation gives
\[
 \dot L_{\mathrm{LP}}=2e^\top X(-Ve+z)+2z^\top NW(-Ve+z).
\]
Since $T$ is diagonal, the sector inequalities $\psi_i(s)(s-\psi_i(s))\ge0$ give
$z^\top T(We-z)\ge0$. Adding twice this quantity bounds $\dot L_{\mathrm{LP}}$ above by the quadratic
form in $(e,z)$ of the matrix in the \emph{Lur'e--Postnikov LMI}
\begin{equation}\label{eq:LPLMI}
 \begin{pmatrix}
 -XV-VX & X-VW^\top N+W^\top T\\
 X-NWV+TW & NW+W^\top N-2T
 \end{pmatrix}\prec0.
\end{equation}
Feasibility of \eqref{eq:LPLMI} therefore implies global exponential stability of $\bar x$. It
certifies stability of that equilibrium rather than contraction, and this function is distinct
from the Fenchel--Bregman divergence.
\bookref{t:KBsharp}{The subcritical contraction statement of Theorem~\ref{t:KBsharp} belongs to
the established theory of coordinatewise slope-restricted Lur'e systems \cite{Davydov}.
Its constructive conclusion is different: with $V,W,L$ fixed, every prescribed finite number of
stable equilibria can be realized by one common scalar response whenever $\bar R_0>1/L$. No conclusion
at $\bar R_0=1/L$ is asserted.}{}
\eeR

For monotone systems and positive linear time-varying systems the converse principle is
\cite[Theorem~4.3, Corollary~4.8]{KawanoBesselinkCao2020}. The next theorem formulates it over a
compact base flow: the base need not be cooperative, $A$ need not be its Jacobian, and only
forward evolution is used.

\beT[Cooperative variational equations admit state-dependent contraction weights exactly when they are uniformly exponentially stable]
[converse contraction weights][Wgt]\label{t:SCM-converse}
Let $g$ be continuous and uniquely integrable near $K$, with continuous
flow $\varphi_t$, and let $A:K\to\mathbb R^{n\times n}$ be continuous and
Metzler. The following are equivalent:
\begin{enumerate}
\item There are $M,\alpha>0$ such that
$\|\Psi(t,y)\|\le Me^{-\alpha t}$ for all $t\ge0$ and $y\in K$.
\item There are $c>0$ and a continuous, uniformly positive weight
$v:K\to\mathbb R^n$ with continuous forward orbital derivative
\[
 \dot v(y)=\left.\frac{d}{dt}v(\varphi_t(y))\right|_{t=0+},
 \qquad \dot v+A^\top v\le-cv\quad\text{on }K.
\]
\end{enumerate}
Under these conditions, every fixed $q\gg0$ gives the weight
\[
 v(y)=\int_0^\infty\Psi(t,y)^\top q\,dt,
 \qquad \dot v+A^\top v=-q.
\]
\eeT
\begin{proof}
The implication (2)$\Rightarrow$(1) follows from
Proposition~\ref{p:VarCert} with $\Sigma=I$ and the minimum and maximum
of the components of $v$ on $K$.
For (1)$\Rightarrow$(2), the displayed integral converges uniformly and
is continuous. Since $A$ is bounded, $\Psi(t,y)=I+O(t)$ uniformly near
$t=0$; positivity gives $v_i(y)\ge\delta q_i/2$ for some $\delta>0$.
The cocycle identity yields, for $s\ge0$,
\[
 v(\varphi_s(y))=
 \Psi(s,y)^{-\top}\left[v(y)-\int_0^s\Psi(t,y)^\top q\,dt\right].
\]
Differentiation at $s=0$ gives the asserted continuous orbital derivative.
Boundedness of $v$ then gives $q\ge cv$ for some $c>0$.
\end{proof}

\beR[Smooth variational equations admit smooth contraction weights]
[][WgtSmooth]
Under the equivalent conditions of Theorem~\ref{t:Wgt}, if $g$ and $A$
extend as $C^1$ fields near $K$, the weight can be chosen $C^1$ near $K$.
Indeed, fix $q\gg0$ and choose $T>0$ so that
$\Psi(T,y)^\top q\le q/2$ on $K$. The finite-horizon weight satisfies
\[
 v_T(y)=\int_0^T\Psi(t,y)^\top q\,dt,\qquad
 D_hv_T+A^\top v_T=\Psi(T,y)^\top q-q\le-q/2.
\]
This is the construction of \cite[Appendix~B, (61)]{KawanoBesselinkCao2020}, there with
$q=\mathbf1$. Finite-time smooth dependence gives $C^1$ regularity; uniform positivity
and boundedness give $D_hv_T+A^\top v_T\le-cv_T$ for some $c>0$.
If $K$ is a rectangular box, $C^1$ tensor Bernstein approximation
provides a polynomial weight: sufficiently close approximants preserve
positivity and the strictly negative residual. Polynomial $g,A$ are
not required for existence, but permit polynomial certificate searches.
\eeR

The next theorem is a compact-domain consequence of the converse principle. Its hypothesis is
global convergence to a Hurwitz equilibrium on a forward invariant set, whereas
\cite[Theorem~4.3]{KawanoBesselinkCao2020} starts from incremental exponential stability on a
forward and backward invariant set.

\beT[Global convergence and contraction weights are equivalent for cooperative systems]
[][EqWgt]
Let $K\subset\mathbb R^n$ be compact, convex, with nonempty interior,
and forward invariant under a $C^2$ vector field $f$ defined near $K$.
Suppose $J(x)=Df(x)$ is Metzler on $K$, $f(e)=0$ for some $e\in K$,
and $J(e)$ is Hurwitz. The following are equivalent:
\begin{enumerate}
\item Every solution starting in $K$ converges to $e$.
\item The state variational matrix $\Psi(t,x)=D_x\varphi_t(x)$ satisfies
$\|\Psi(t,x)\|\le Me^{-\alpha t}$ for some $M,\alpha>0$, all $t\ge0$,
and all $x\in K$.
\item There are $c>0$ and a $C^1$ weight $v\gg0$ near $K$ such that
\[
 (\partial_xv)f+J^\top v\le-cv\qquad\text{on }K.
\]
\end{enumerate}
In this case $e$ is globally exponentially stable relative to $K$.
If $K$ is a rectangular box, $v$ may be chosen polynomial.
\eeT
\begin{proof}
Theorem~\ref{t:Wgt} and Remark~\ref{r:WgtSmooth} give (2)$\Leftrightarrow$(3)
and the polynomial conclusion. Convexity gives
\[
 \varphi_t(x)-e
 =\int_0^1\Psi(t,e+\theta(x-e))(x-e)\,d\theta,
\]
so (2) implies exponential convergence and hence (1).

Assume (1). Since $J(e)$ is Metzler and Hurwitz,
\[
 w=(-J(e))^{-\top}\mathbf1\gg0,\qquad P_w=\operatorname{Diag}(w).
\]
For $\mu_1(H)=\max_j(H_{jj}+\sum_{i\ne j}|H_{ij}|)$, continuity gives
$\lambda>0$ and a neighbourhood $U$ of $e$ such that
$\mu_1(P_wJ(x)P_w^{-1})\le-\lambda$ on $U\cap K$.
Local asymptotic stability gives a relatively open forward-invariant
neighbourhood $V$ of $e$ in $K$, contained in $U$.
Every $x\in K$ enters $V$ at some time $t_x$; continuity gives a
relatively open neighbourhood of $x$ entering $V$ at the same time.
A finite subcover and forward invariance give $T<\infty$ such that
$\varphi_t(K)\subset V$ for all $t\ge T$.
With $L_w=\sup_{x\in K}\|P_wJ(x)P_w^{-1}\|_1$, the initial bound
$e^{L_wt}$ and subsequent decay at rate $\lambda$ yield
\[
 \|P_w\Psi(t,x)P_w^{-1}\|_1
 \le e^{(L_w+\lambda)T}e^{-\lambda t}
 \qquad(t\ge0,\ x\in K).
\]
Equivalence of norms proves (2).
\end{proof}

In the setting of Theorem~\ref{t:EqWgt}, the weight $v$ of item (3) defines the intrinsic distance
\[
 d_v(x,y)=\inf_\gamma\int_0^1v(\gamma(\theta))^\top|\gamma'(\theta)|\,d\theta,
\]
over $C^1$ paths in $K$ from $x$ to $y$. Item (3) and the Metzler property give
$d_v(\varphi_t(x),\varphi_t(y))\le e^{-ct}d_v(x,y)$, so $L(x)=d_v(x,e)$ is an equilibrium Lyapunov
function, possibly nonsmooth, with $L(\varphi_t(x))\le e^{-ct}L(x)$.

\emph{\shortalt{Transverse application}{Transverse and compound applications}.}
For a resident flow, taking $A=J^\perp$ in
Proposition~\ref{p:VarCert} or Theorem~\ref{t:Wgt} certifies the transverse
variational equation without requiring a cooperative resident field.
Theorem~\ref{t:EqWgt}, in contrast, concerns the full Jacobian; its
convergence conclusion cannot be inferred from transverse decay alone.

\beP[Composite Lyapunov functions combine resident storage with transverse decay][][ComLya]
Write the equations as $\dot z=f_z(z,y)$, $\dot y=f_y(z,y)$ with $z\ge0$ and $f_z(0,y)=0$, and let
$y^*$ be a resident equilibrium. On a compact forward-invariant set $K\subseteq\{z\ge0\}$ containing
$(0,y^*)$, suppose there are $\ell\gg0$, $a,c>0$, $b\ge0$, and a continuous resident storage
$U\ge0$, vanishing only at $y^*$ and differentiable near the resident projection of $K$, such that
\[
 \ell^\top f_z(z,y)\le-c\ell^\top z,
\]
\[
 \nabla U(y)^\top f_y(0,y)\le-a\|y-y^*\|^2,
 \qquad
 \nabla U(y)^\top\bigl(f_y(z,y)-f_y(0,y)\bigr)\le b\ell^\top z
 \quad ((z,y)\in K).
\]
Then, for any $\kappa>b/c$, the function $H(z,y)=U(y)+\kappa\ell^\top z$ satisfies
\[
 \dot H\le-a\|y-y^*\|^2-(\kappa c-b)\ell^\top z,
\]
and $(0,y^*)$ is globally asymptotically stable relative to $K$.
\eeP
\begin{proof}
Adding the three inequalities gives the estimate for $\dot H$. Since $z\ge0$ and $\ell\gg0$, $H\ge0$ on
$K$ and vanishes only at $(0,y^*)$, and the right-hand side is negative elsewhere. Thus $H$ is a strict
Lyapunov function on the compact forward-invariant set $K$.
\end{proof}
This is a standard composite-Lyapunov argument. The first inequality concerns the full nonlinear
equations on $K$, not merely their transverse linearization on the siphon face.
\bookref{s:FenBre}{The linear term is the zero-reference limit of Subsection~\ref{s:FenBre}.}{}
\bookref{p:ComBas}{Proposition~\ref{p:ComBas} gives an application to a threshold network.}{}

% Compound and Orlando discussion: unused by the short paper.
\ifshort\else
For $\dot x=f(x)$ of class $C^1$, take $g=f$ and
$A=J^{[k]}$, where $J=Df$ and $J^{[k]}$ is the additive compound on
$\bigwedge^k\mathbb R^n$, of dimension $\binom nk$.
If $\Psi$ is the state variational matrix, differentiation of its minors gives
\[
 \frac{d}{dt}\Psi^{(k)}=J^{[k]}\Psi^{(k)},\qquad \Psi^{(k)}(0)=I.
\]
Proposition~\ref{p:VarCert} therefore bounds infinitesimal
$k$-dimensional volumes. Linearity of $J\mapsto J^{[k]}$ carries a
convex enclosure by $J_1,\ldots,J_q$ to the enclosure by
$J_1^{[k]},\ldots,J_q^{[k]}$.

\emph{Orlando's identity ($n\ge2$).}
For $p(\lambda)=\det(\lambda I-J)=\lambda^n+a_1\lambda^{n-1}+\cdots+a_n$,
let $\Delta_{n-1}(p)=\det(a_{2j-i})_{i,j=1}^{n-1}$, with $a_0=1$
and $a_\ell=0$ outside $0\le\ell\le n$. Since the eigenvalues of
$J^{[2]}$ are $\lambda_i+\lambda_j$, $i<j$, Orlando's identity reads
\begin{equation}\label{HG:orlandoeq}
 \Delta_{n-1}(p)=\det(-J^{[2]})
 =(-1)^{\binom n2}\prod_{i<j}(\lambda_i+\lambda_j).
\end{equation}
At a Hopf point a pair sums to zero. Thus a contraction certificate
for $J^{[2]}$ excludes a Hopf point at any equilibrium where it holds.
\fi

% Inactive legacy result retained for source compatibility.
\iffalse
\beP[Autonomous convergence by a constant compound norm][][CompConv]
Let $\dot x=f(x)$ be an autonomous system with $f$ of class $C^1$
on an open neighbourhood of a compact, convex, forward invariant
set $K\subset\mathbb R^n$, $n\ge2$, with nonempty interior in the
chosen reduced coordinates. Put $J(x)=Df(x)$ and
$N_2=\binom n2$. Let $J^{[2]}$ be the additive compound acting on
$\bigwedge^2\mathbb R^n$. Suppose there are a constant invertible
$\Theta\in\mathbb R^{N_2\times N_2}$, a vector norm with induced
logarithmic norm $\mu$, and a constant $c>0$ such that
\[
 \mu\bigl(\Theta J^{[2]}(x)\Theta^{-1}\bigr)\le-c
 \qquad(x\in K).
\]
Then $K$ contains no nonconstant periodic orbit and every trajectory
starting in $K$ approaches the set
$\mathcal E=\{x\in K:f(x)=0\}$.
If $\mathcal E=\{e\}$, every such trajectory converges to $e$.
If $e$ is also locally asymptotically stable relative to $K$, it is
globally asymptotically stable relative to $K$.
\eeP
\begin{proof}[Proof by the autonomous convergence theorem]
On compound coordinates use the fixed norm
$\|\zeta\|_*:=\|\Theta\zeta\|$. Its logarithmic norm satisfies
$\mu_*(B)=\mu(\Theta B\Theta^{-1})$.
The hypothesis therefore gives precisely the constant-norm condition of
\cite[Theorem~12]{WuKanevskiyMargaliot2022} on the compact convex
invariant set $K$. That theorem gives convergence to the equilibrium
set. A nonconstant periodic orbit cannot approach that closed set,
since its compact image is disjoint from it. The uniqueness assertion
and the final stability assertion follow immediately.
Only a neighbourhood of $K$ is needed: a smooth cutoff equal to one
near $K$ extends the vector field to all of $\mathbb R^n$ without
changing its trajectories or Jacobian on $K$.
\end{proof}
\fi
